\chapter{Series RLC Circuits: Resonance, Damping, and Q}
\label{ch:rlc}
The series RLC circuit is central to tuned receivers, oscillators, matching
networks, traps, crystal filters, and every selective RF stage. With two energy
stores---a capacitor and an inductor---it is the electrical embodiment of the
\emph{second-order system} of \cref{ch:linsys}, and everything that chapter said
about \(\omega_0\), \(\zeta\), damping regimes, pole geometry, and resonant peaks
applies here directly. Its dual, the parallel RLC circuit, follows in
\cref{ch:rlcpar}.

\section{The Circuit}
A source drives \(R\), \(L\), and \(C\) in series (\cref{fig:rlc-schematic}). Two
independent states carry the stored energy: the capacitor voltage \(v_C\) and the
inductor current \(i\). We take the output across the capacitor,
\(v_o=v_C\), which is the low-pass response; taking it across \(R\) gives a
band-pass and across \(L\) a high-pass, all from the same circuit.

\begin{figure}[htbp]
\centering
\begin{circuitikz}[scale=1]
\draw (0,0) to[sV,l=\(u(t)\)] (0,2)
      to[R,l=\(R\)] (2,2)
      to[L,l=\(L\),i>=\(i(t)\)] (4,2)
      to[C,l=\(C\),v=\(v_C\)] (4,0) -- (0,0);
\draw (4,2) to[short,-o] (5.1,2) node[right]{\(+\)};
\draw (4,0) to[short,-o] (5.1,0) node[right]{\(-\)};
\node at (5.5,1) {\(v_o\)};
\end{circuitikz}
\caption{Series RLC circuit, with the output taken across the capacitor
(\(v_o=v_C\)). The two states are the capacitor voltage \(v_C\) and the inductor
current \(i\), so the system is second order.}
\label{fig:rlc-schematic}
\end{figure}

\section{From Schematic to Differential Equation}
Start from the physics, one relation per element. KVL around the single loop says
the source equals the sum of the three element voltages, and the capacitor law
relates its current to its rate of voltage change:
\[
u(t)=R\,i+L\frac{\dd i}{\dd t}+v_C,
\qquad
i=C\frac{\dd v_C}{\dd t}.
\]
Written in the two state variables, this is a \emph{coupled system of two
first-order differential equations}---one per energy store:
\[
\underbrace{\frac{\dd v_C}{\dd t}=\frac{1}{C}\,i}_{\text{capacitor law}},
\qquad
\underbrace{\frac{\dd i}{\dd t}=\frac{1}{L}\bigl(u-R\,i-v_C\bigr)}_{\text{KVL}} .
\]
This pair is the most faithful statement of the circuit: two stores, two states,
two equations. To reach the familiar single second-order equation, substitute.
Differentiate the first equation, \(\dd^2v_C/\dd t^2=(1/C)\,\dd i/\dd t\), insert
the second, and replace \(i=C\,\dd v_C/\dd t\):
\[
\frac{\dd^2 v_C}{\dd t^2}
=\frac{1}{LC}\Bigl(u-RC\frac{\dd v_C}{\dd t}-v_C\Bigr)
\qquad\Longrightarrow\qquad
\boxed{\,LC\frac{\dd^2 v_C}{\dd t^2}+RC\frac{\dd v_C}{\dd t}+v_C=u(t).\,}
\]
The two views are equivalent; the coupled pair generalizes better, while the single
second-order equation matches the standard form of \cref{ch:linsys} directly.

\section{State-Space Form}
The coupled pair above \emph{is} the state-space model; writing it in matrix form
with \(x_1=v_C\) and \(x_2=i\),
\[
\dot{\state}=
\underbrace{\mat{0 & 1/C\\ -1/L & -R/L}}_{A}\state
+\underbrace{\mat{0\\ 1/L}}_{B}u,
\qquad
y=v_o=v_C=\underbrace{\mat{1 & 0}}_{C_{\text{out}}}\state .
\]
Two states, two poles---the eigenvalues of \(A\). Computing them explicitly,
\[
\det(sI-A)
=\det\!\mat{s & -1/C\\ 1/L & s+R/L}
=s\Bigl(s+\frac{R}{L}\Bigr)+\frac{1}{LC}
=s^2+\frac{R}{L}s+\frac{1}{LC},
\]
so the eigenvalues are
\[
s_{1,2}=-\frac{R}{2L}\pm\sqrt{\left(\frac{R}{2L}\right)^{\!2}-\frac{1}{LC}} .
\]
Keep these two numbers in view: the next section reaches them again by a completely
different route.

\section{Matching the Generic Second-Order System}
Divide the ODE by \(LC\) to reach the standard form
\(\ddot v_C+2\zeta\omega_0\dot v_C+\omega_0^2 v_C=\omega_0^2 u\), giving the
identifications
\[
\omega_0=\frac{1}{\sqrt{LC}},
\qquad
\alpha=\frac{R}{2L},
\qquad
\zeta=\frac{\alpha}{\omega_0}=\frac{R}{2}\sqrt{\frac{C}{L}},
\qquad
Q=\frac{1}{2\zeta}=\frac{\omega_0 L}{R}.
\]
The poles follow from the standard form as
\(s=-\zeta\omega_0\pm\jj\omega_0\sqrt{1-\zeta^2}=-\alpha\pm\jj\omega_d\) with
\(\omega_d=\sqrt{\omega_0^2-\alpha^2}\).

\begin{controlsbox}
\textbf{The poles are the eigenvalues.} These must be the same two numbers found
from \(\det(sI-A)=0\) above, and they are. Substituting
\(\alpha=R/2L\) and \(\omega_0^2=1/LC\) into the eigenvalue expression,
\[
s_{1,2}=-\frac{R}{2L}\pm\sqrt{\left(\frac{R}{2L}\right)^{2}-\frac{1}{LC}}
=-\alpha\pm\sqrt{\alpha^2-\omega_0^2}
=-\alpha\pm\jj\sqrt{\omega_0^2-\alpha^2}
=-\alpha\pm\jj\omega_d,
\]
where the last step assumes the underdamped case \(\alpha<\omega_0\) (which is what
makes the square root imaginary). Eigenvalues of the state matrix, roots of the
characteristic polynomial, and poles of the transfer function are three names for
one pair of numbers---exactly as \cref{ch:linsys} promised.
\end{controlsbox}
Now the second-order results of \cref{ch:linsys} transfer wholesale: the step
responses of \cref{fig:step2} (over/critical/under-damped), the pole geometry of
\cref{fig:poles2} (on the \(\omega_0\) circle, \(\cos\theta=\zeta\)), and the
resonant peak of \cref{fig:bode2}. A tuned RF circuit is simply a very
lightly damped (high-\(Q\)) instance.

\section{Series Resonance}
The series impedance is
\[
Z=R+\jj\left(\omega L-\frac{1}{\omega C}\right).
\]
At \(\omega=\omega_0\) the reactances cancel, \(Z=R\) is purely resistive, and the
current is maximum for a fixed applied voltage. The reactive voltages across \(L\)
and \(C\) individually can then exceed the source voltage by a factor of \(Q\).

\section{Q and Bandwidth}
\label{sec:halfpower}
The quality factor measures sharpness. For the series circuit,
\[
Q_s=\frac{\omega_0 L}{R}=\frac{1}{\omega_0 RC},
\]
and the half-power bandwidth is \(BW=f_0/Q\). That is usually quoted with an
\(\approx\), as though it were a narrowband estimate. It is not: for this circuit it
is an identity, and deriving it also reveals where the band edges actually sit.

Factor \(R\) out of the impedance. Using the first form of \(Q_s\) for the inductive
term (\(\omega L/R=Q\omega/\omega_0\)) and the second for the capacitive one
(\(1/\omega CR=Q\omega_0/\omega\)),
\[
Z=R\left[1+\jj Q\,D(\omega)\right],
\qquad
D(\omega)\equiv\frac{\omega}{\omega_0}-\frac{\omega_0}{\omega},
\qquad\text{so}\qquad
|Z|^2=R^2\left[1+Q^2D(\omega)^2\right].
\]
Every appearance of frequency is now inside the single \emph{detuning} function
\(D\), which vanishes at resonance. The current is \(I=V/Z\), so the power delivered
to \(R\) is \(P=|I|^2R=V^2R/|Z|^2\), largest at \(\omega_0\) where \(|Z|=R\). Half
that power means \(|Z|^2=2R^2\), which is simply
\[
Q\,D(\omega)=\pm1 .
\]
Two exact consequences follow, and neither needs \(Q\) to be large.

\textbf{The edges straddle resonance geometrically.} Substitute \(\omega_0^2/\omega\)
for \(\omega\):
\(D(\omega_0^2/\omega)=\omega_0/\omega-\omega/\omega_0=-D(\omega)\). Reflecting a
frequency through \(\omega_0\) in this reciprocal sense flips the sign of the
detuning, so it maps the \(+1\) solution onto the \(-1\) solution: the two half-power
frequencies are precisely such a pair. Hence
\[
\omega_1\omega_2=\omega_0^2,
\qquad\text{that is}\qquad
f_0=\sqrt{f_1f_2}\;.
\]
Resonance is the \emph{geometric} mean of the band edges, not the midpoint between
them.

\textbf{The width is exactly \(\omega_0/Q\).} That result licenses writing the edges
as \(\omega_2=r\omega_0\) and \(\omega_1=\omega_0/r\) for a single \(r>1\). The upper
edge satisfies \(QD=+1\), which is \(Q(r-1/r)=1\). But the bandwidth is that same
combination:
\[
BW=\omega_2-\omega_1=\omega_0\left(r-\frac{1}{r}\right)=\frac{\omega_0}{Q},
\qquad\text{so}\qquad
BW=\frac{f_0}{Q}
\]
with no approximation anywhere in the argument. Solving the same quadratic gives the
edges themselves,
\(r=\tfrac{1}{2Q}+\sqrt{1+1/4Q^2}\).

The two results also say how wrong it is to centre a band arithmetically. Since
\(r+1/r=\sqrt{(r-1/r)^2+4}\),
\[
\frac{\omega_1+\omega_2}{2}=\frac{\omega_0}{2}\left(r+\frac{1}{r}\right)
=\omega_0\sqrt{1+\frac{1}{4Q^2}}
\approx\omega_0\left(1+\frac{1}{8Q^2}\right),
\]
so the midpoint of the two \(-\SI{3}{\decibel}\) frequencies sits above resonance by
a fraction \(1/8Q^2\)---\SI{2.5}{ppm} at the \(Q\approx224\) of the box below,
\SI{68}{ppm} at the \(Q\approx43\) of this chapter's worked example, and still only
\SI{0.4}{\percent} at \(Q=5.7\). That is why splitting the difference on a bench
sweep never bites, and why the exact statement went unnoticed.

On the Bode magnitude (\cref{fig:bode2}) higher \(Q\) therefore means a taller,
narrower peak---poles pressed against the \(\jj\omega\) axis, and it confirms
\cref{sec:polegeometry}'s reading of the pole diagram: \(\omega_0/Q=2\alpha\) is
exactly twice the distance from that axis, so ``distance from the imaginary axis is
half the bandwidth'' is an equality too. \Cref{ch:rlcpar} shows that the parallel circuit's \(Q\)
is the \emph{reciprocal} combination, and why.

\begin{mistakebox}
Two things about \(BW=f_0/Q\) do \emph{not} generalize, and neither is the equals
sign.

\textbf{The \(\approx\) you have seen elsewhere belongs to the output.} The equality
holds for the \emph{current}---equivalently for the voltage across \(R\), the bandpass
output---because that is the response whose magnitude is governed by \(|Z|\) alone.
Take the output across \(C\) instead and the peak's \(-\SI{3}{\decibel}\) width
matches \(\omega_0/Q\) only as \(Q\) grows: it is \SI{0.05}{\percent} wide of it at
\(Q=22\) but \SI{6}{\percent} off at \(Q=2.2\), where the edges are no longer
geometrically centred either. So the familiar hedge is not about high-\(Q\)
sloppiness---it is about which variable you called the output.

\textbf{The formula for \(Q\) is topology-dependent, and there are three of them.}
Do not carry one topology's \(Q\) into another, and keep the three straight:
\emph{loaded} \(Q\) includes source and load coupling; \emph{unloaded} \(Q\)
describes the resonator's own internal losses; \emph{component} \(Q\) describes the
loss of one inductor or capacitor. They answer different questions.
\end{mistakebox}

\begin{workedbox}[: the voltage magnification is real]
A series circuit with \(L=\SI{20}{\micro\henry}\), \(C=\SI{100}{\pico\farad}\) and
\(R=\SI{2}{\ohm}\) is driven by \SI{10}{\volt} RMS at resonance. Then
\[
f_0=\frac{1}{2\pi\sqrt{LC}}\approx\SI{3.56}{\mega\hertz},
\qquad
\omega_0L\approx\SI{447}{\ohm},
\qquad
Q_s=\frac{447}{2}\approx224 .
\]
At resonance \(Z=R=\SI{2}{\ohm}\), so the current is \(10/2=\SI{5}{\ampere}\) and
the voltage across the inductor alone is
\[
V_L=I\,\omega_0L=(5)(447)\approx\SI{2.2}{\kilo\volt},
\]
which is \(Q\) times the \SI{10}{\volt} source---and the capacitor holds an equal
and opposite \SI{2.2}{\kilo\volt}, which is why their sum contributes nothing and
the source only ever sees \SI{2}{\ohm}. This is not a curiosity: it is why the
capacitors in a high-\(Q\) tuned circuit or antenna tuner must be rated for
kilovolts even when the transmitter output is modest.
\end{workedbox}

\section{Root Locus: Damping Set by R}
Fix \(L\) and \(C\) and sweep \(R\). Since \(\zeta=\tfrac{R}{2}\sqrt{C/L}\), the
poles trace the canonical second-order locus (\cref{fig:rlc-rootlocus}): at
\(R=0\) they sit on the \(\jj\omega\) axis at \(\pm\jj\omega_0\) (undamped,
infinite \(Q\)); as \(R\) grows they arc along the \(\omega_0\) circle into the
left half-plane, meet at \(-\omega_0\) when \(R=2\sqrt{L/C}\) (critical), then
split along the real axis (overdamped). Selectivity, ringing, and damping are the
one moving pole pair seen three ways. This is the general second-order locus of
\cref{fig:rootlocus2} with \(R\) as the sweep parameter.

\begin{figure}[htbp]
\centering
\begin{tikzpicture}[>=Latex,scale=1.15]
\draw[->,black!70] (-3.6,0) -- (0.9,0) node[right]{\(\sigma\) (real)};
\draw[->,black!70] (0,-2.6) -- (0,2.6) node[above]{\(\jj\omega\) (imag)};
\draw[GuideLine,dashed] (0,0) circle (2);
\node[GuideLine] at (1.55,1.75) {\(|s|=\omega_0\)};
\filldraw[GuideBlue] (0,2) circle (2.4pt);
\filldraw[GuideBlue] (0,-2) circle (2.4pt);
\node[GuideBlue,right] at (0.08,2.15) {\small \(R=0\)};
\filldraw[GuideGreen] (-1.732,1.0) circle (2.4pt);
\filldraw[GuideGreen] (-1.732,-1.0) circle (2.4pt);
\filldraw[GuideAmber] (-2,0) circle (2.4pt);
\node[GuideAmber,below left] at (-2.0,-0.08) {\small critical};
\filldraw[black!70] (-1.25,0) circle (2.4pt);
\filldraw[black!70] (-3.2,0) circle (2.4pt);
\draw[->,GuideGreen,thick] (0,1.98) to[bend right=18] (-1.7,1.05);
\draw[->,GuideGreen,thick] (0,-1.98) to[bend left=18] (-1.7,-1.05);
\draw[->,black!70,thick] (-1.9,0.05) to[bend right=25] (-1.3,0.02);
\draw[->,black!70,thick] (-2.1,0.05) to[bend left=25] (-3.15,0.02);
\draw[black!55] (0,0) -- (-1.732,1.0);
\node at (-0.55,0.55) {\small \(\theta\)};
\node[GuideGreen] at (-2.55,1.35) {\small increasing \(R\)};
\end{tikzpicture}
\caption{Root locus of the series RLC poles as \(R\) increases with \(L,C\)
fixed: from \(\pm\jj\omega_0\) on the axis (undamped), along the circle of radius
\(\omega_0\), meeting at \(-\omega_0\) at critical damping (\(R=2\sqrt{L/C}\)), then
splitting along the real axis (overdamped). Here \(\cos\theta=\zeta\).}
\label{fig:rlc-rootlocus}
\end{figure}

\section{The Same Circuit as a Second-Order Filter}
\label{sec:lc-lowpass}
Nothing new is needed to use this circuit as a filter---the transfer function is
already in hand. Laplace-transforming the boxed ODE (zero initial conditions),
\[
\bigl(LC\,s^2+RC\,s+1\bigr)V_o(s)=U(s)
\qquad\Longrightarrow\qquad
G(s)=\frac{V_o}{U}=\frac{1}{LC\,s^2+RC\,s+1}
=\frac{\omega_0^2}{s^2+\dfrac{\omega_0}{Q}\,s+\omega_0^2},
\]
using \(\omega_0/Q=2\zeta\omega_0=R/L\). This is exactly the standard second-order
low-pass of \cref{ch:linsys}: flat at low frequency, a resonant peak of height
\(\approx Q\) near \(\omega_0\) when \(Q>1/\sqrt2\), then falling at
\(-\SI{40}{\decibel}\) per decade because there are two poles.

In practice the series \(R\) is rarely a deliberate resistor. It is the coil's loss
plus the source and load resistance the section works between---which is why the
\emph{same} \(L\) and \(C\) give a peaky response when lightly loaded and a smooth
one when properly terminated.

\begin{controlsbox}
\textbf{The two design knobs fall right out.} From the first-principles model,
\(\omega_0=1/\sqrt{LC}\) sets the cutoff frequency and
\(Q=(1/R)\sqrt{L/C}\) sets the peaking---how close the pole pair sits to the
\(\jj\omega\) axis. Scaling \(L\) and \(C\) together slides the whole response
along the frequency axis without changing its shape; changing the resistance
(loss, or the source/load termination) changes only the damping. Every passive
filter design is some version of choosing these two, times the number of poles
(\cref{ch:filters}).
\end{controlsbox}

\section{Energy Balance}
With \(\mathcal{H}=\tfrac12 Cv_C^2+\tfrac12 Li^2\), the series energy balance is
\[
\frac{\dd\mathcal{H}}{\dd t}=u(t)\,i(t)-R\,i^2(t).
\]
\begin{energybox}
The inductor and capacitor trade energy back and forth at \(\omega_0\); the resistor
is the only drain. High \(Q\) means little is lost per cycle, so the exchange rings
for many cycles---the time-domain face of a sharp resonance, and the reason
\(Q=2\pi\times\)(stored)/(lost per cycle) is the definition that survives across
topologies (\cref{ch:rlcpar}).
\end{energybox}

\begin{controlsbox}
High \(Q\) \(\Leftrightarrow\) low \(\zeta\) \(\Leftrightarrow\) poles near the
\(\jj\omega\) axis \(\Leftrightarrow\) narrow bandwidth and long ringing. The tuned
circuit the exam calls ``selective'' is what a controls course calls ``lightly
damped.''
\end{controlsbox}

\begin{exambox}
Series-resonance questions reward separating the three quantities that do not
interact. \(f_0=1/(2\pi\sqrt{LC})\) depends on \(L\) and \(C\) \emph{only}---any
\(R\) in the problem statement is a distractor for this part. \(Q=\omega_0L/R\) and
\(BW=f_0/Q\) then depend on \(R\), and \(R\) alone moves them. So a question
that changes \(R\) changes selectivity and bandwidth but never the resonant
frequency; a question that changes \(L\) or \(C\) moves \(f_0\), and moves \(Q\)
too unless it scales both. At resonance \(Z=R\), the current is maximum, and the
individual reactive voltages are \(Q\) times the source.
\end{exambox}

\section{Worked Example}
Let \(L=\SI{10}{\micro\henry}\), \(C=\SI{220}{\pico\farad}\), \(R=\SI{5}{\ohm}\).
Then
\[
f_0=\frac{1}{2\pi\sqrt{LC}}\approx\SI{3.39}{\mega\hertz},
\qquad
Q_s=\frac{\omega_0 L}{R}\approx 42.6,
\qquad
\zeta=\frac{1}{2Q}\approx0.0117.
\]
The poles sit at \(s=-\alpha\pm\jj\omega_d\approx-2.5\times10^5\pm\jj\,2.13\times10^7\ \si{\radian\per\second}\):
the real part is about \(1/85\) of the imaginary part, so the poles are pressed
against the \(\jj\omega\) axis---a very high-\(Q\), sharply peaked circuit that
rings for roughly \(Q\approx43\) cycles. The half-power bandwidth is
\(BW=f_0/Q\approx\SI{80}{\kilo\hertz}\), and by \cref{sec:halfpower} the two
\(-\SI{3}{\decibel}\) frequencies are \SI{3.35}{\mega\hertz} and
\SI{3.43}{\mega\hertz}---whose \emph{product} is \(f_0^2\).

\Cref{fig:rlc-worked} plots this particular circuit on real axes. The magnitude
peaks about \(\SI{33}{\decibel}\) (a factor of \(43\)) above the passband at
\SI{3.39}{\mega\hertz} and then falls at \(-\SI{40}{\decibel}\) per decade; the
phase swings the full \(0^\circ\to-180^\circ\) that two poles must supply, passing
\(-90^\circ\) exactly at \(f_0\). On the locus, the design point sits a hair off the
\(\jj\omega\) axis---which is what ``\(Q=43\)'' looks like geometrically.

\begin{figure}[htbp]
\centering
\includegraphics[width=\textwidth]{rlc_series_worked}
\caption{The worked example, \(L=\SI{10}{\micro\henry}\),
\(C=\SI{220}{\pico\farad}\), \(R=\SI{5}{\ohm}\). \emph{Left:} Bode magnitude and
phase of \(v_C/u\) on a real frequency axis, with the \(Q\approx43\) resonant peak
at \(f_0=\SI{3.39}{\mega\hertz}\). \emph{Right:} root locus as \(R\) is swept with
\(L,C\) fixed. The design point (\(R=\SI{5}{\ohm}\)) is pressed against the
\(\jj\omega\) axis; the poles would meet on the real axis only at the critical value
\(R=2\sqrt{L/C}\approx\SI{426}{\ohm}\), some 85 times larger.}
\label{fig:rlc-worked}
\end{figure}

\subsection{The Same Circuit as an Impedance Phasor}
It is worth converting these same numbers into the phasor language of
\cref{ch:ac}, since that is how a bench instrument reports them. Just above
resonance, at \(f=1.1f_0\),
\[
X_L=\omega L\approx\SI{235}{\ohm},
\qquad
X_C=\frac{1}{\omega C}\approx\SI{194}{\ohm},
\qquad
X=X_L-X_C\approx\SI{41}{\ohm},
\]
so the impedance phasor is
\[
Z=R+\jj X\approx 5+\jj41\ \si{\ohm}
=\SI{41}{\ohm}\angle 83^\circ .
\]
Two things are worth noticing (\cref{fig:rlc-phasor}). First, the individual
reactances are huge---hundreds of ohms---yet they nearly cancel, leaving a net
\SI{41}{\ohm}; \emph{resonance is a cancellation, not an absence}. Second, sweeping
frequency moves \(Z\) along a \emph{vertical} line in the complex plane, because
\(R\) is constant and only the reactance changes. That line crosses the real axis at
\(Z=R\)---the one frequency where the circuit looks purely resistive, and where
\(|Z|\) is smallest and the current therefore largest.

\begin{figure}[htbp]
\centering
\includegraphics[width=\textwidth]{rlc_series_phasor}
\caption{The same worked example seen as an impedance phasor. \emph{Left:}
\(Z=5+\jj41\ \si{\ohm}=\SI{41}{\ohm}\angle83^\circ\) at \(f=1.1f_0\)---rectangular
form on the axes, polar form along the arrow. \emph{Right:} the locus of \(Z(f)\)
from \(0.5f_0\) to \(2f_0\): a vertical line at \(\Re\{Z\}=R\), inductive above
\(f_0\) and capacitive below, crossing the real axis at resonance.}
\label{fig:rlc-phasor}
\end{figure}

